\section{Experiments}

\subsection{Setup}

\paragraph{Scenes.}
We evaluate on 14 scenes: 7 raw UAV scenes (\texttt{raw\_self}, \texttt{raw001}--\texttt{raw006}) and 7 official aligned scenes (\texttt{ms\_bud}, \texttt{ms\_fruit}, \texttt{ms\_garden}, \texttt{ms\_golf}, \texttt{ms\_lake}, \texttt{ms\_solar}, \texttt{ms\_tree}). Raw scenes test rectification, QA, reconstruction, products, and geometry diagnostics. Official scenes remove rectification confounds and are used for external common-mask comparison.

\paragraph{Methods.}
We report two UMGS variants: UMGS-I, an independent per-band transfer pipeline, and UMGS-J, a joint band-bank pipeline with one authoritative multispectral checkpoint. MS-Splatting is the external baseline on official aligned scenes. Scratch-MS and Unfrozen-support are representative single-band upper-bound ablations; they are not promoted as product models because they do not preserve the shared Gaussian support required by model-level multispectral fusion.

\paragraph{Metrics.}
Image quality: PSNR/SSIM/LPIPS. Spectral means average G/R/RE/NIR. Index fidelity: MAE/RMSE/PSNR/SSIM for NDVI, GNDVI, NDRE. Geometry: relative-depth AbsRel and agreement at 1\%, 5\%, 10\%. Cost: non-invasive GPU traces and exported Gaussian counts. Official external comparison uses a common-mask protocol; UMGS shims are resolution-aligned comparison views only.

\subsection{Main Results on 14 Scenes}

Table~\ref{tab:main14} shows that UMGS-I and UMGS-J remain near-identical on aggregate quality and geometry. On all 14 scenes, both reach 0.820 spectral SSIM and 0.806 index SSIM; UMGS-I has a small edge on index RMSE (0.084 vs.\ 0.085). Depth aggregates are identical because both variants share the same locked support.

This result supports a stable interpretation: UMGS-I remains the conservative primary line, while UMGS-J is a valid joint-checkpoint branch with numerically comparable performance.

\begin{table*}[!t]
    \caption{Main UMGS results over 14 scenes. RGB metrics for UMGS-J use the shared UMGS-I RGB anchor. Spectral metrics average G/R/RE/NIR. Index metrics average NDVI/GNDVI/NDRE. Depth uses the relative reference-depth protocol and averages G/R/RE/NIR.}
    \label{tab:main14}
    \centering
    \scriptsize
    \setlength{\tabcolsep}{3pt}
    \begin{tabular}{llcccccccccc}
        \toprule
        Method & Split & RGB PSNR & RGB SSIM & MS PSNR & MS SSIM & MS LPIPS & Idx. RMSE & Idx. SSIM & AbsRel & Ag@5\% & Ag@10\% \\
        \midrule
        UMGS-I & Raw 7 & 24.39 & 0.840 & 27.40 & 0.850 & 0.136 & 0.065 & 0.869 & 0.019 & 0.916 & 0.988 \\
        UMGS-I & Official 7 & 23.30 & 0.741 & 26.07 & 0.791 & 0.248 & 0.103 & 0.743 & 0.117 & 0.697 & 0.750 \\
        UMGS-I & All 14 & 23.85 & 0.791 & 26.74 & 0.820 & 0.192 & 0.084 & 0.806 & 0.068 & 0.807 & 0.869 \\
        \midrule
        UMGS-J & Raw 7 & 24.39 & 0.840 & 27.38 & 0.850 & 0.136 & 0.067 & 0.869 & 0.019 & 0.916 & 0.988 \\
        UMGS-J & Official 7 & 23.30 & 0.741 & 26.07 & 0.791 & 0.249 & 0.103 & 0.743 & 0.117 & 0.697 & 0.750 \\
        UMGS-J & All 14 & 23.85 & 0.791 & 26.72 & 0.820 & 0.192 & 0.085 & 0.806 & 0.068 & 0.807 & 0.869 \\
        \bottomrule
    \end{tabular}
\end{table*}


\subsection{Index-Focused Analysis}

Table~\ref{tab:index-detail} reports NDVI/GNDVI/NDRE separately. Differences between UMGS-I and UMGS-J are small on NDVI and GNDVI, with a slightly larger NDRE gap favoring UMGS-I in RMSE/PSNR. The joint-bank design is therefore viable but not consistently superior.

\begin{table}[!t]
    \caption{UAV-focused 10-scene spectral-index fidelity by vegetation index.}
    \label{tab:index-detail}
    \centering
    \small
    \begin{tabular}{llcccc}
        \toprule
        Method & Index & MAE & RMSE & PSNR & SSIM \\
        \midrule
        UMGS-I & NDVI & 0.058 & 0.083 & 28.49 & 0.847 \\
        UMGS-I & GNDVI & 0.059 & 0.083 & 28.65 & 0.823 \\
        UMGS-I & NDRE & 0.041 & 0.059 & 31.53 & 0.864 \\
        \midrule
        UMGS-J & NDVI & 0.058 & 0.083 & 28.49 & 0.847 \\
        UMGS-J & GNDVI & 0.059 & 0.083 & 28.66 & 0.823 \\
        UMGS-J & NDRE & 0.043 & 0.063 & 31.31 & 0.863 \\
        \bottomrule
    \end{tabular}
\end{table}


\subsection{Official-Aligned External Comparison}

Table~\ref{tab:mms-official7} summarizes the 7-scene common-mask comparison against MS-Splatting. On the 7-scene average, UMGS-I-shim and UMGS-J-shim are higher in SSIM and lower in LPIPS, spectral RMSE, and SAM. Scene-level behavior in Table~\ref{tab:mms-scenes} is mixed: \texttt{ms\_lake} and \texttt{ms\_golf} favor MS-Splatting on selected entries, while other scenes favor our geometry-locked shims. We report this block as competitive and scene dependent rather than universally dominant.

\begin{table}[!t]
    \caption{Official-aligned 7-scene common-mask comparison against MS-Splatting. Lower is better for LPIPS, 4-band RMSE, and SAM. UMGS shims are resolution-aligned comparison views, not separate training variants.}
    \label{tab:mms-official7}
    \centering
    \small
    \begin{tabular}{lccccc}
        \toprule
        Method & PSNR & SSIM & LPIPS & 4-band RMSE & SAM$^\circ$ \\
        \midrule
        UMGS-I-shim & 26.56 & 0.820 & 0.231 & 0.0544 & 5.75 \\
        UMGS-J-shim & 26.55 & 0.820 & 0.231 & 0.0544 & 5.75 \\
        MS-Splatting & 26.03 & 0.803 & 0.272 & 0.0571 & 6.23 \\
        \bottomrule
    \end{tabular}
\end{table}

\begin{table*}[!t]
    \caption{Scene-level aligned aerial common-mask comparison for UMGS-I$^\dagger$ and MS-Splatting. PSNR/SSIM/LPIPS average G/R/RE/NIR; 4-band RMSE and SAM evaluate the spectral vector.}
    \label{tab:mms-scenes}
    \centering
    \scriptsize
    \setlength{\tabcolsep}{3pt}
    \begin{tabular}{lcccccccccc}
        \toprule
        Scene & UMGS PSNR & MMS PSNR & UMGS SSIM & MMS SSIM & UMGS LPIPS & MMS LPIPS & UMGS RMSE & MMS RMSE & UMGS SAM & MMS SAM \\
        \midrule
        \texttt{ms\_golf} & 31.33 & 32.12 & 0.928 & 0.935 & 0.176 & 0.234 & 0.0300 & 0.0273 & 2.77 & 3.11 \\
        \texttt{ms\_lake} & 22.02 & 22.79 & 0.740 & 0.731 & 0.350 & 0.374 & 0.0887 & 0.0817 & 8.84 & 7.88 \\
        \texttt{ms\_solar} & 27.72 & 26.21 & 0.893 & 0.893 & 0.117 & 0.123 & 0.0448 & 0.0532 & 4.52 & 5.16 \\
        \bottomrule
    \end{tabular}
\end{table*}


For raw unaligned data, Table~\ref{tab:retained-self} is an applicability diagnostic under the external retained-subset path. It is not the main raw benchmark because it does not represent the full raw protocol used by UMGS-I/UMGS-J.

\begin{table}[!t]
    \caption{Raw UMGS-Field retained-subset diagnostic. This is an applicability check for the external pipeline under its retained subset, not the main raw full-scene comparison.}
    \label{tab:retained-self}
    \centering
    \scriptsize
    \begin{tabular}{lccccc}
        \toprule
        Method & PSNR & SSIM & LPIPS & 4-band RMSE & SAM$^\circ$ \\
        \midrule
        UMGS-I (retained) & 28.35 & 0.834 & 0.203 & 0.0502 & 3.98 \\
        UMGS-J (retained) & 28.33 & 0.834 & 0.204 & 0.0503 & 4.00 \\
        MS-Splatting (default) & 31.07 & 0.933 & 0.107 & 0.0379 & 3.95 \\
        \bottomrule
    \end{tabular}
\end{table}



\subsection{Representative Ablations}

Table~\ref{tab:ablations3} compares the locked UMGS variants with two upper-bound ablations on \texttt{raw\_self}, \texttt{raw002}, and \texttt{ms\_lake}. The unfrozen branch improves selected image metrics, but it is not a valid product model: on \texttt{ms\_lake}, product construction fails the shared-support audit with an xyz mismatch of $6.84\times10^{-2}$. Scratch-MS diverges more strongly, including Gaussian-count mismatch across bands. These failures occur at the exact stage where multispectral products are assembled, showing that per-band quality gains alone are insufficient if support consistency is lost.

\begin{table*}[!t]
    \caption{Representative 3-scene ablation summary on \texttt{raw\_self}, \texttt{raw002}, and \texttt{ms\_lake}. Spectral metrics average G/R/RE/NIR; index metrics average NDVI/GNDVI/NDRE. Scratch-MS and Unfrozen-support are single-band upper-bound ablations, not valid product models, because their shared-support audits fail.}
    \label{tab:ablations3}
    \centering
    \scriptsize
    \setlength{\tabcolsep}{3pt}
    \begin{tabular}{lcccccc}
        \toprule
        Method & MS PSNR & MS SSIM & MS LPIPS & Idx. RMSE & Shared & Products \\
        \midrule
        UMGS-I & 25.35 & 0.786 & 0.228 & 0.100 & yes & yes \\
        UMGS-J & 25.33 & 0.786 & 0.228 & 0.100 & yes & yes \\
        Scratch-MS & 25.88 & 0.780 & 0.289 & 0.116 & no & invalid \\
        Unfrozen-support & 26.27 & 0.794 & 0.218 & 0.105 & no & invalid \\
        \bottomrule
    \end{tabular}
\end{table*}


\subsection{Cost}

Table~\ref{tab:cost} reports coarse but traceable cost indicators. UMGS-J uses higher recorded peak VRAM but lower average recorded VRAM in the aggregated traces; both UMGS variants retain identical Gaussian counts. The external baseline has much smaller recorded VRAM on official adapter runs, but that value should not be interpreted as full-pipeline raw-scene cost.

\begin{table}[!t]
    \caption{Coarse recorded cost indicators on the UAV-focused suite. Values are aggregated from non-invasive GPU traces and exported PLY vertex counts; they are intended as a first cost table, not a microbenchmark.}
    \label{tab:cost}
    \centering
    \scriptsize
    \begin{tabular}{lcccc}
        \toprule
        Method & Peak VRAM GB & Avg. VRAM GB & Peak util. \% & Gaussians M \\
        \midrule
        UMGS-I & 27.45 & 3.99 & 99.0 & 4.31 \\
        UMGS-J & 35.04 & 1.50 & 99.0 & 4.31 \\
        \bottomrule
    \end{tabular}
\end{table}

